\documentclass[../main.tex]{subfiles}

\graphicspath{{../image/}{../figure/}}

\begin{document}

\chapter{随机选择法}\label{chapter:ch13-随机选择法}

本章介绍著名的随机选择法（random choice method），它为构造严格双曲守恒律 Cauchy 问题的全局定义、可容许 $BV$ 解提供了最早且至今仍十分有效的格式；其中初始数据的总变差很小。该解是一列不抹平激波的近似解的极限。在第九章中详细讨论过的 Riemann 问题的解，充当构造 Cauchy 问题近似解的建筑砌块。力图保持激波的陡峭，可能与算法相容性的要求相冲突。该方法的``随机性''特征被用来达成这样一种微妙的平衡：在保证相容性的同时，不抹平传播激波锋面的陡峭性。波追踪（wave tracing）这一装置将在此仅作简要讨论，它以描绘全局波型为代价，使算法变为确定性的。

我们将针对特征族要么真非线性、要么线性退化的系统，给出随机选择法的细致介绍。更一般系统的情形涉及实质性的技术复杂性，此处仅作扼要提及。

在第十六章中，我们将看到如何把该算法改造为处理带非齐次项与源项的平衡律系统。

\subfile{Chapter-13/section-01.tex}
\subfile{Chapter-13/section-02.tex}
\subfile{Chapter-13/section-03.tex}
\subfile{Chapter-13/section-04.tex}
\subfile{Chapter-13/section-05.tex}
\subfile{Chapter-13/section-06.tex}
\subfile{Chapter-13/section-07.tex}
\subfile{Chapter-13/section-08.tex}
\subfile{Chapter-13/section-09.tex}

\end{document}